\documentclass[UTF8,11pt]{ctexart}
\usepackage{amsmath,amssymb,bm}
\usepackage{booktabs}
\usepackage{graphicx}
\usepackage{hyperref}
\usepackage{cleveref}
\usepackage{xcolor}
\usepackage{enumitem}
\usepackage{microtype}
\usepackage{geometry}
\usepackage{caption}
\usepackage{subcaption}
\usepackage[backend=biber,style=ieee,sorting=none]{biblatex}

\geometry{a4paper,margin=25mm}
\graphicspath{{../system_overview/assets/final/}{../system_overview/assets/selected/}}
\addbibresource{references.bib}
\hypersetup{colorlinks=true,linkcolor=blue!55!black,citecolor=blue!55!black,urlcolor=blue!55!black}

\newcommand{\systemname}{\textsc{PRE-MAP-VLN}}
\newcommand{\pose}{\bm{q}}
\newcommand{\todo}[1]{\textcolor{red!70!black}{[TODO: #1]}}

\newcommand{\overviewfigure}{%
  \begin{figure*}[t]
    \centering
    \IfFileExists{../system_overview/assets/final/system_overview.pdf}{%
      \includegraphics[width=\textwidth]{system_overview.pdf}%
    }{%
      \fbox{\parbox[c][45mm][c]{0.94\textwidth}{\centering
      System Overview placeholder\\
      Export the final figure to\\
      \texttt{paper/system\_overview/assets/final/system\_overview.pdf}}}%
    }
    \caption{Overview of \systemname. The system first constructs a reusable 3D semantic map, then jointly reasons about language constraints, executable observation poses, and collision-free flight paths. Online observations trigger conditional replanning or semantic recovery when needed.}
    \label{fig:overview}
  \end{figure*}%
}


\title{面向室内无人机长时程语言任务的三维语义地图与语义--几何联合规划}
\author{作者姓名（待填写）}
\date{\today}

\begin{document}
\maketitle

\begin{abstract}
现有视觉语言导航方法通常以到达目标邻域作为成功标准，难以同时处理长指令中的部分顺序、条件分支、空间关系以及任务相关的最终观察姿态。本文提出 \systemname，一套面向室内无人机的预建图长时程任务执行系统。系统首先通过三维自主探索建立占据地图、分层点云、物体三维框、房间区域和跨层通道；随后将自然语言解析为带约束的任务图，根据任务语义在三维地图中生成可执行观察位姿，并联合优化任务顺序、终止位姿与三维飞行路径。执行阶段利用开放词汇检测进行在线验证；当历史位置失效时，系统结合大模型语义先验、历史场景实例、三维可达代价和负观测结果开展主动恢复搜索。当前系统已在 Habitat 多层住宅场景中完成端到端实现，包括完整三维 A*、碰撞检查、B 样条轨迹、跨层执行和 RViz 可视化。实验将评估任务成功率、路径长度、规划时间、观察有效性与恢复搜索效率。
\end{abstract}

\noindent\textbf{关键词：}视觉语言导航；室内无人机；三维语义地图；长时程任务；观察位姿；主动搜索

\overviewfigure

\section{引言}
服务机器人接收到的自然语言往往包含多个目标，例如检查不同房间中的物体、满足必要的先后关系，并根据当前观测决定是否执行后续分支。局部反应式导航缺少整条任务的全局空间规划能力；直接使用全局地图的方法又常把目标简化为二维点，并将“靠近目标”视为任务完成。对于具有固定相机安装方式的室内无人机，最终位置、飞行高度和航向会直接影响目标是否进入视野，因此目标邻域并不等价于可执行的终止状态。

本文研究具有预建三维语义地图的室内长时程语言任务。核心问题是：在保留语言中的必要约束时，如何决定任务顺序、为每个任务选择合适的三维观察位姿，并生成安全且总代价较低的连续飞行轨迹；当地图中的目标先验与当前观测不一致时，又如何主动恢复。

本文的主要贡献暂定如下：
\begin{itemize}[leftmargin=*]
  \item 构建一条面向多层室内无人机的三维预探索与语义建图链路，联合表达占据体素、物体实例、房间区域和跨层通道。
  \item 将长指令解析为具有部分顺序、条件分支、空间关系和失败策略的任务图，而不是机械遵循原句顺序。
  \item 根据任务语义和锚点三维框生成带位置与航向的可执行观察位姿，并联合优化任务顺序、观察位姿和三维路径。
  \item 提出地图先验失效后的语义引导主动搜索机制，将大模型语义候选、历史实例、三维可达性和在线负观测统一用于恢复决策。
\end{itemize}

\section{相关工作}
本节将从四方面整理相关研究：无全局地图的视觉语言导航、基于语义地图或场景图的导航、任务相关观察位姿规划，以及面向无人机的探索与轨迹规划。正式版本需补充 SpatialNav、VLMaps、SayNav、CondVLN、FALCON、OccuSG 与 Boxer 等工作的准确引文和定量差异。\todo{补齐文献与对比表。}

\section{问题定义}
环境由三维几何地图 $\mathcal{M}_g$ 和语义地图 $\mathcal{M}_s$ 表示。语言指令被解析为任务图 $\mathcal{G}_T=(\mathcal{V},\mathcal{E})$，其中节点表示检查、寻找或确认任务，边表示必要先后关系和条件依赖。任务 $i$ 的候选终止位姿集合为
\begin{equation}
G_i=\{\pose_i^1,\pose_i^2,\ldots,\pose_i^{K_i}\},\qquad
\pose=(x,y,z,\psi),
\end{equation}
其中 $\psi$ 为无人机航向。系统需要选择满足任务图约束的任务顺序 $\pi$ 与每个任务的候选位姿 $k_i$，使总执行代价最小：
\begin{equation}
\min_{\pi,\{k_i\}}\sum_j C_{\mathrm{motion}}
(\pose_{\pi_j}^{k_{\pi_j}},\pose_{\pi_{j+1}}^{k_{\pi_{j+1}}};\mathcal{M}_g)
+\lambda C_{\mathrm{view}}(\pose_{\pi_j}^{k_{\pi_j}}),
\end{equation}
同时保证三维可达、碰撞安全、任务可见性与语言约束。条件观测会改变有效任务集合，因此系统需在线重规划。

\section{系统方法}
\subsection{三维自主探索与语义地图构建}
FALCON 使用在线体素地图执行完整三维探索，并通过三维 A*、碰撞检查和 B 样条轨迹驱动 Habitat 中的无人机。探索结果经 Boxer 提取物体三维框，再由墙体结构、楼层聚类、跨层通道重建和 OccuSG 生成房间区域。房间标签由场景内物体清单辅助推断。该阶段不读取真值楼层作为在线规划依据，所得地图可在后续任务中复用。

\subsection{语言任务图}
大语言模型将用户指令转换为结构化任务，输出目标类别、房间或楼层范围、空间关系、任务目的、未找到策略、条件分支和必要前置关系。场景校验器只验证输出是否能与当前场景图对齐，不替代模型决定可重排的任务顺序。用户明确指定的位置始终作为初始搜索范围；只有失败策略允许时才扩大搜索范围。

\subsection{任务相关观察位姿}
系统将语义关系归一化为支撑区域、下方区域、周边区域和实例本体等几何类型。对锚点水平尺寸 $w,d$，基础采样半径为
\begin{equation}
r=r_{\mathrm{box}}+d_{\mathrm{obs}},\qquad
r_{\mathrm{box}}=\frac{1}{2}\sqrt{w^2+d^2}.
\end{equation}
候选高度由观察区域高度、水平距离和相机垂直视场约束；航向始终朝向观察区域中心。之后删除碰撞、不可达、视锥外、严重遮挡和覆盖率过低的候选。对于同一锚点，多视角按方位角连续访问，减少停顿和不必要的航向往复。

\subsection{语义--几何联合规划}
候选层使用共享占据语义的粗三维 A* 估算真实可达代价，不可达边返回无穷大而非欧氏距离。规划器在满足任务图偏序和条件约束的前提下，联合选择任务顺序与代表观察位姿。最终执行仍由权威三维 A*、碰撞验证与 B 样条优化完成。路径代价缓存绑定几何地图版本，避免地图更新后复用失效路径。

\subsection{在线验证与语义恢复}
无人机到达候选位姿后运行开放词汇目标检测，并将 RGB-D 检测投影到全局坐标。若一个位置的有效多视角均未找到目标，则判定该历史位置失效。系统向大模型提供原始指令、失败位置和场景物体清单，获得可能的房间与功能区域；再将其映射为实际三维锚点。下一搜索位置综合语义先验、三维路径代价、观察质量和负观测历史进行选择。目标被发现后，更新语义地图并重新优化剩余任务。

\section{实验设计}
\subsection{平台与场景}
实验平台为 Habitat 室内仿真器，主体场景为多层住宅 00337-CFVBbU9Rsyb，并计划在七个 HM3D 场景上验证泛化性。无人机具有三维平移与航向控制能力，相机不独立优化俯仰角。\todo{补充传感器参数、地图分辨率、安全距离和硬件。}

\subsection{任务与指标}
实验包含两类任务：（1）带部分顺序、条件分支和空间关系的多目标长指令；（2）指定位置失败后触发的语义恢复搜索。主要指标包括任务完成率、语言约束满足率、最终目标可见率、路径长度、规划时间、恢复搜索访问房间数、误检率和碰撞次数。

\subsection{基线与消融}
拟比较按语言顺序执行、仅按欧氏距离排序、二维候选代价、无观察位姿优化、无条件重规划以及无语义恢复等变体。\todo{根据最终实验范围确定基线，避免报告尚未完成的结果。}

\section{当前结果与讨论}
当前 00337 场景已完成多层三维探索、物体框提取、墙体与房间分割、跨层通道重建、长指令解析、三维候选规划、在线 OWLv2 验证、条件分支和目标丢失恢复。RViz 可同步展示任务状态、活动楼层点云、已执行轨迹、候选视锥、目标框及检测帧。正式结果表应只采用可复现实验输出，并严格区分成功完成、超时、规划失败和人工终止。\todo{导入最终定量结果与可视化。}

\section{局限性}
当前语义地图依赖预探索质量；开放词汇检测仍可能产生类别混淆；大模型输出需要结构化校验；Habitat 场景中的高处几何与真实无人机传感器存在差异。此外，房间语义、动态物体更新和跨场景阈值仍需更系统的验证。

\section{结论}
本文给出一套将三维语义地图、长指令约束、任务相关观察位姿和安全飞行轨迹统一起来的室内无人机任务系统。该方法把导航目标从“靠近一个点”扩展为“到达能够完成任务的三维观察状态”，并在地图先验失效时通过语义与几何联合恢复继续执行。后续工作将集中于多场景实验、定量消融和真实无人机验证。

\printbibliography
\end{document}
